% mode: plain
% A name not found exactly is looked for without regard to case, on the disk of the directory where it would be.  The directory's
% entries are kept between misses, so a file the job writes itself must be seen by the very next look, whatever the clock of the
% directory says.  An \input of a file whose name differs by case from what is asked for, before and after it exists.
\catcode`\@=11
\immediate\openin1=casefoldwritten.tex
\ifeof1 \immediate\write16{no such file yet}\else \immediate\write16{found}\fi
\immediate\closein1
\immediate\openin1=CASEFOLDWRITTEN.tex
\ifeof1 \immediate\write16{no such file yet}\else \immediate\write16{found}\fi
\immediate\closein1
\immediate\openout2=CaseFoldWritten.tex
\immediate\write2{\string\def\string\fromfile{read back}}
\immediate\closeout2
\immediate\openin1=casefoldwritten.tex
\ifeof1 \immediate\write16{no such file}\else \immediate\write16{found under another case}\fi
\immediate\closein1
\input CASEFOLDWRITTEN.TEX
\immediate\write16{\fromfile}
\immediate\openin1=CASEFOLDWRITTEN2.tex
\ifeof1 \immediate\write16{no such file}\else \immediate\write16{found}\fi
\immediate\closein1
\immediate\openout2=casefoldwritten2.TEX
\immediate\write2{\string\def\string\fromfile{read back again}}
\immediate\closeout2
\immediate\openin1=CASEFOLDWRITTEN2.tex
\ifeof1 \immediate\write16{no such file}\else \immediate\write16{found under another case}\fi
\immediate\closein1
\input CaseFoldWritten2.tex
\immediate\write16{\fromfile}
\end
